\paragraph{Conjugate-coordinate class: the contraction by the ambient metric and the five-matrix presentation (conditional; bounded no-go and conditional route).}
Keep the Poisson limit, the leaf \(\mathbb{CP}^{1,2}\) of \(\mathfrak{so}(4,2)\) with \(\eta_6=\mathrm{diag}(-1,1,1,1,1,-1)\), the notation \(y^a=M^{a5}\), \(y^4=M^{45}\), \(Q^{ab}=(y^4)^2\eta^{ab}+y^ay^b\) of the preceding paragraph, and \(T^a=M^{a4}\), \(D^a=\{T^a,\cdot\}\). On the orbit \((\eta_6M)^2=-\ell^2\mathbf 1\) (checked at \(200\) random points to \(10^{-12}\)), whence \(y\cdot y=-\ell^2\) and \(T^ay_a=0\). The base metric induced from the ambient space is \(g=\eta+y_\mu y_\nu/(y^4)^2\) with inverse \(g^{-1}=\eta+yy/\ell^2\); it is positive definite, and the contraction of the fields \(D^a\) by \(Q^{-1}\) has base symbol exactly \(g^{-1}\) (checked numerically). The kinetic tensor of the four matrices \(T^a\) is \(K=(y^4)^2\eta=(y^4)^2\bigl(g^{-1}-yy/\ell^2\bigr)\). The contraction by \(Q^{-1}\), equivalently the term \((y_aD^a)^2/\ell^2\) added to \(\eta_{ab}D^aD^b\), therefore replaces the Lorentzian tensor \(K\) by the Euclidean metric \(g\) of the hyperboloid; it is not a Wick rotation of \(\eta\).

\emph{No supply in the four-matrix class.} The Hessian of the quartic term plus the bilinear term at \(T^a\), \(L(A)_a=-\Box A_a+D_a(D_bA^b)-2s\mathcal J_{ab}A^b+\mu^2A_a\), contains the matrix indices only through the constant form \(\eta_{ab}\), and the ghost operator is \(-\Box\). A gauge-fixing term adds \(G^\dagger G\ge0\) to the gauge-invariant Hessian, so the number of negative modes of the gauge-fixed operator does not exceed that of the gauge-invariant one, and a regulator that is a function of operators does not change the signs of the eigenvalues. The field \(y_aD^a=-T^a\{y_a,\cdot\}\) is not in the Lie algebra generated by the \(D^a\), which is \(\mathfrak{so}(4,1)'\); the functions \(y_a\) are not polynomials in the \(T^a\) of bounded degree. Hence in the class of the preceding paragraph none of the action, the bilinear term, the background gauge, the ghosts and the \(\Box_X\) regulator produces the term \((y_aD^a)^2\). The full Casimir \(\tfrac12\eta_{AC}\eta_{BD}J^{AB}J^{CD}\) has base symbol \(-2\ell^2g^{-1}\) (checked numerically), but it enters neither the Hessian nor the ghost operator. Extensions of the matrix set by \(y^a\) with a positive-norm contraction (eight matrices \(T^a,y^a\); seven matrices \(T^a,y^i\)) have rank-six horizontal span, but the equation of motion then needs different bilinear coefficients for different matrices (checked by exact Poisson algebra), so they require further inputs.

\emph{Five-matrix presentation.} Take \(X^A=M^{Ac_0}\), \(A\ne c_0\), with the restriction of \(\eta_6\) as contraction (Definition~\ref{def:matrix-coords-ext}, \(D=5\)). The orbit relation gives \(y\cdot y=\eta_{c_0c_0}\ell^2\) for \(y^A=M^{Ac_0}\). At \(25\) random points for each \(c_0\) the horizontal rank is \(4\) and the centre rank \(2\); the induced contraction \(h=c^{T}\eta c\) on the horizontal space is positive definite for \(c_0\in\{0,5\}\) and has signature \((2,2)\) for \(c_0\in\{1,\dots,4\}\); a Lorentzian signature does not occur. For \(X^A=M^{A5}\), \(A=0,\dots,4\) (stabiliser \(\mathfrak{so}(4,1)\)): \(y_AD^A=0\), the base symbol of \(\eta_{AB}D^AD^B\) is \(\ell^2g^{-1}\) (checked to \(10^{-12}\)), and at \(300\) random points and \(91\) directions of \(\zeta\), for both signs of the bracket, the vertical flux in an \(h\)-orthonormal frame is self-dual or anti-self-dual to \(10^{-11}\), so \(f_1=f_2\) identically. The bilinear coefficient required by stationarity is \(\mu^2=+4\) (exact, both signs of \(s\)), the Zeeman term of the gauge-fixed vector operator has the same modulus as in the four-matrix class, and in the polynomial Poisson algebra, exactly,
\[
H_V=C_{\mathfrak{so}(4,2)}-C_{G}+C_{\Sigma}+\mu^2,\qquad \Box=C_{\mathfrak{so}(4,1)}-C_{\mathfrak{so}(4,2)},
\]
where \(C_G\) is the Casimir of the total \(\mathfrak{so}(4,1)\) action on vector-valued functions and \(C_\Sigma\) that of the vector index. By the Weyl formulae the \(\mathfrak{so}(4,1)\) component \((k+1,k)\) of the degree-\(k\) vector polynomials has \(C_{\mathfrak{so}(4,2)}-C_G+C_\Sigma=0\) for every \(k\), so it is an eigenspace of \(H_V\) with eigenvalue \(\mu^2\) (verified directly for \(k=1\), dimension \(35\)). At the leading symbol the vector operator then has the lowest Landau level at zero for the four horizontal components, a zero-mode family with density \(f^2/(2\pi^2)\) in \(\zeta\), independent of the proper time; the normal component has no Zeeman term, and the bosonic density per \(\mathrm d^2\zeta/(2\pi)^2\) is \(f^2(16\pi^2)^{-1}\bigl[8+(3+n_{\mathrm{ex}})/\sinh^2(sf)\bigr]\). The \(\zeta\)-integral of the second term converges; that of the first diverges and is not regulated by the proper-time cutoff. The fibre of the leaf is compact, so \(\zeta\) is bounded by the fibre harmonic number of the quantum space; the divergence is a feature of the Poisson limit.

\emph{Compact check at finite size.} For the compact form \(\mathfrak{so}(6)\) with \(\delta\) throughout, on the \(N\)-th symmetric power of the fundamental representation, the exact Hessian of the quartic term plus \(\tfrac12m^2\operatorname{Tr}X^2\) (finite matrices, \(N=1,\dots,4\)) satisfies \(H+G^\dagger G=C_6-C_5^{\mathrm{tot}}+C_5^{\mathrm{vec}}+m^2\) for \(X^A=M^{A6}\), \(m^2=-4\). Its lowest eigenvalue is \(-4=m^2\) for \(N\le4\), with the family \((k+1,k)\), \(k\le N\), of total multiplicity \(\sum_{k\le N}(2k+1)(2k+3)(2k+5)/3\) and \(14\) further modes at \(-2\); for the four-matrix background \(X^a=M^{a5}\), \(m^2=-3\), it is \(-(N+3)\), growing with \(N\). The sign of \(m^2\) is \(-4\) on the sphere and \(+4\) on the hyperboloid.

\emph{Consequences.} (i) The term \((y_aD^a)^2\) is supplied by the matrix data only if the matrices are the position functions \(X^A=M^{A5}\); the price is that the emergent horizontal metric is then the Euclidean metric \(g\), and the conformally flat Lorentzian \(K\) of the four-matrix class is not obtained, so the preceding level-one and stress results do not carry over. (ii) For the hyperboloid, the vector zero modes of the leading symbol are lifted by \(+\mu^2\) only if the vanishing of the Casimir combination, which is exact on polynomial vector fields and on the compact form, continues to the normalisable sector; this continuation is not shown. (iii) Neither the sign nor the size of \(w\) follows: the weight-two coefficient and its variation under \(A\) are not computed, and the passage from a Euclidean coefficient to the Lorentzian one is a Wick-rotation hypothesis tied to Problem~E.

This is a statement for the Poisson limit, the leading symbol, equiregular points, the bosonic content and the single bilinear term; the self-duality of the flux is verified numerically and not proved, and the finite-size spectrum is for the compact form. It does not treat fermions, Moyal terms for the hyperboloid, or the Lorentzian emergent metric. This changes no ledger label; Problem~F remains open.
